\documentclass[12pt]{article}
\usepackage[utf8]{inputenc}
\usepackage[T1]{fontenc}
\usepackage{hyperref}
\title{Bwanabwana Language Resources for K-12 Curriculum Development}
\author{Compiled by Claude Sonnet 5 (claude-sonnet-5) for Andr\'as Kornai}
\date{}
\begin{document}
\maketitle

\section{Introduction}
Bwanabwana (also Tubetube; ISO 639-3 tte) is a Western Oceanic Austronesian language spoken by about 2,400 people on small islands in the Louisiade Archipelago, Milne Bay Province, Papua New Guinea (Tubetube, Ole, Anagusa, Kitai, Wale and Kwalaiwa). A local newspaper article describes two main dialects, Tubetube and Wale.

\section{Primary Educational Resources}
None found.

\section{Linguistic and Reference Materials}
\begin{itemize}
\item Wikipedia: \url{https://en.wikipedia.org/wiki/Bwanabwana_language}
\item Omniglot: \url{https://www.omniglot.com/writing/bwanabwana.htm}
\end{itemize}

\section{Religious and Community Resources}
\begin{itemize}
\item Joshua Project reports that the New Testament and portions of Scripture exist, and that a full Bible is still needed. \url{https://joshuaproject.net/index.php/people_groups/15611}
\item The National (PNG newspaper) article on Bible translation into the language, with SIL work on Tubetube island. \url{https://www.thenational.com.pg/?p=376750}
\item eBible page titled ``Buki Tabu Beyabeyana'', which appeared in the search results. I have not checked that it is the Bwanabwana text, and its code (tteo) differs from tte. \url{https://ebible.org/tteo/FRT01.htm}
\end{itemize}

\section{Digital and Audio Resources}
No audio or video resources were found.

\section{Curriculum Development Context}
No school use was found, which matches the reference description. The Scripture translation is the largest body of native-language text found. The related languages Dobu, Tawala, Misima, Duau and Suau share many words with Bwanabwana, according to The National article. No primers or dictionaries were found.
\end{document}